\section{Conclusions and Future Work}
\label{conc_section}

This paper introduced GreenMeasurer, a lightweight and extensible framework designed to facilitate non-intrusive measurement of energy consumption and runtime performance of software systems on Linux and macOS. The framework reads hardware energy counters through backends (Intel RAPL, NVIDIA NVML and, on macOS, \textit{powermetrics}), selects the CPU and GPU backends independently, reports unsupported metrics explicitly, and incorporates temperature calibration, over-time monitoring and, on Intel CPUs, power capping. Experimental results, including the execution of a Fibonacci algorithm in C, illustrate GreenMeasurer's capability to produce reproducible energy measurements, as well as insights into the impact of power limits on energy consumption and performance metrics. Measurements are hardware-domain measurements, not per-process attributions.

Several limitations remain. The macOS backends are implemented but not yet validated on hardware, the AMD backend is not validated, and the API has no explicit ABI version, so the core and backends must be built from compatible sources. Future work includes validated AMD support, a comparison with CodeCarbon and EnergiBridge (Section~\ref{sec:eval_plan}), a comparison against an external power meter, and support for virtualized environments and further components such as I/O and network.
